% Generated draft: requires mathematical and source review.
\documentclass[11pt]{article}
\usepackage[margin=25mm]{geometry}
\usepackage{fontspec}
\setmainfont{Latin Modern Roman}
\usepackage{amsmath,amssymb,mathtools,unicode-math}
\setmathfont{Latin Modern Math}
\usepackage{enumitem,needspace,xurl,hyperref}
\setlength{\parindent}{0pt}
\setlength{\parskip}{5pt}
\setlength{\emergencystretch}{3em}
\newcommand{\N}{\mathbb{N}}
\newcommand{\Z}{\mathbb{Z}}
\newcommand{\Q}{\mathbb{Q}}
\newcommand{\R}{\mathbb{R}}
\newcommand{\C}{\mathbb{C}}
\newcommand{\F}{\mathbb{F}}
\newcommand{\A}{\mathbb{A}}
\newcommand{\PP}{\mathbb{P}}
\newcommand{\E}{\mathbb{E}}
\newcommand{\eps}{\varepsilon}
\newcommand{\dd}{\,\mathrm{d}}
\begin{document}
\section*{UnsolvedMath ID: 30005279}
\textbf{Source problem code:} catalogue entry
\subsection{Definitions and mathematical statement}
This record combines six questions numbered 9--14 in the supplied extract. Several lack hypotheses or a complete target; the qualifications below are part of this statement-only draft.

\begin{enumerate}
\item \textbf{Degree-five superirreducibility (source question 9).} The supplied definition calls a polynomial \(f\in\Z[x]\) of degree at least two \emph{\(2\)-superirreducible} if \(f(g(x))\) is irreducible whenever \(g\in\Z[x]\) has degree at most two. The question is whether there exists such an \(f\) with \(\deg f=5\). As written, the degree condition includes nonzero constant polynomials, so it includes constant compositions and leaves the intended irreducibility convention unresolved. Restricting instead to \(1\leq\deg g\leq2\) would change the supplied definition and requires confirmation. The source also does not specify whether irreducibility is taken in \(\Z[x]\) or \(\Q[x]\), or impose primitivity on \(f\).

The supplied background attributes to Schinzel (1967) the assertion that, for \(f\in\Z[x]\) with \(\deg f=d\geq3\), there is a polynomial \(g\) of degree \(d-1\) such that \(f(g(x))\) is reducible. The extract does not explicitly repeat the coefficient domain of this \(g\) or provide a separate citation. This background assertion is retained as supplied attribution.

\item \textbf{Power subgroups of unit groups (source question 10).} For a positive modulus \(n\), write \(U_n=(\Z/n\Z)^\times\). For an integer \(d\), define its power subgroup by
\[
H_{n,d}=\{z^d:z\in U_n\},
\]
with the multiplication inherited from \(U_n\). The supplied source asserts that every finite abelian group \(G\) admits a representation \(G\cong H_{n,d}\), and asks for sharp upper bounds for the parameters \(n\) and \(d\) in terms of \(|G|\). The assertion of representability is supplied background, not a result established here. A precise optimization problem requires a specified admissible range for \(d\), conventions for the trivial group and modulus \(1\), and a choice between separate bounds and bounds achieved by the same representation. These choices are absent from the extract.

\item \textbf{An identity involving a norm form (source question 11).} For the intended, but unspecified, field extension \(K/\Q\) and elements \(\omega_1,\ldots,\omega_5\in K\), the displayed target is to find \(f,g\in\Z[x_1,\ldots,x_5]\), homogeneous of degrees \(3\) and \(2\), respectively, satisfying the polynomial identity
\[
f(x_1,\ldots,x_5)^2+g(x_1,\ldots,x_5)^3
=N_{K/\Q}\!\left(\sum_{i=1}^{5}x_i\omega_i\right).
\]
Here \(N_{K/\Q}\) denotes the field norm when \(K/\Q\) is a finite extension. Neither \(K\), its degree, nor the role or quantification of the \(\omega_i\) is supplied. In particular, five variables do not establish that these elements are a basis. The compatibility of the norm form with the degree-six left-hand side and with integral coefficients cannot be checked from the record. No particular extension or choice of elements is assumed here.

\item \textbf{Double covering by prime residue classes (source question 12).} For a positive integer \(N\), choose a finite set \(P\) of distinct primes and one residue \(a_p\in\Z/p\Z\) for each \(p\in P\). Define
\[
c(m)=\#\{p\in P:m\equiv a_p\pmod p\}.
\]
The required double covering means
\[
c(m)\geq2\qquad\text{for every }m\in\{1,\ldots,N\}.
\]
The two supplied questions ask whether such a covering can be obtained (a) using only primes \(p\leq N\), and (b) with
\[
\sum_{p\in P}\frac1p\leq K.
\]
The record does not specify whether these questions concern every \(N\), all sufficiently large \(N\), or another range. It also leaves unspecified the quantifiers on \(K\), its permitted dependence on \(N\), and whether part (b) retains the restriction \(p\leq N\). The letter \(K\) here is a bound, unrelated to the field in question 11.

\item \textbf{Residues of consecutive primes (source question 13).} The original extract supplies the following background assertion, attributed to Daniel Shiu (2000): for a positive modulus \(q\) and an integer \(a\) with \(\gcd(a,q)=1\), there are arbitrarily long strings of consecutive primes all congruent to \(a\pmod q\). Writing \(p_j\) for the \(j\)-th positive prime, this says that for every positive integer \(k\) there exists \(j\geq1\) such that
\[
p_{j+i}\equiv a\pmod q\qquad(0\leq i<k).
\]
This is supplied background, not an open question asserted here. The reconstructed statement requests arbitrary prescribed patterns of consecutive prime residues, but the original extract gives no corresponding question. For specified \(q\), \(k\), and residue classes \(a_0,\ldots,a_{k-1}\), an occurrence of such a pattern would mean
\[
p_{j+i}\equiv a_i\pmod q\qquad(0\leq i<k)
\]
for some \(j\geq1\). The admissible patterns and the intended quantifiers, including whether infinitely many occurrences are required, are missing; this notation does not assert that every pattern is admissible or occurs.

\item \textbf{Quaternionic analytic methods (source question 14).} Develop an analogue of the circle method, Weyl sums, or Poisson summation for quaternions. This is a research direction rather than a specified theorem: the quaternion algebra, coefficient domain, discrete set or order, functions or sums, analytic estimates, and intended applications are not supplied. A more precise mathematical target cannot be recovered from the record alone.
\end{enumerate}
\subsection{Short English statement}
The record collects six questions about polynomial composition, power subgroups, norm forms, double coverings, consecutive prime residues, and quaternionic analytic methods. Several formulations need additional hypotheses or missing source context. The supplied classification is open, but neither that classification nor the mathematical claims have been independently verified here.
\subsection{Sources}
\begin{itemize}
\item Local UnsolvedMath registry, \texttt{data/unsolvedmath/problems.json}, permanent ID \(30005279\), title ``Six Problems from an Analytic Number Theory Session,'' exact source code \texttt{catalogue entry}. The supplied record has status \texttt{open} and statement status \texttt{corrected\_verified}; its claim of checking against the official report is supplied provenance, not independent verification by this draft.
\item Supplied report citation: \emph{Analytic Number Theory (2023)}, Oberwolfach Reports, pp.~2895--2960, DOI \texttt{10.4171/owr/2022/50}. Supplied source link: \url{https://doi.org/10.4171/owr/2022/50}. The record gives proposed year \(2022\), while the citation names \(2023\); these supplied date labels are retained.
\item Supplied literature source: ``Oberwolfach Report 50,'' \url{https://ems.press/content/serial-article-files/46986}. The literature assessment dated \texttt{2026-08-22} classifies the combined record as open, citing unresolved covering, prime-pattern, and quaternionic Weyl-sum questions. This is supplied dated curation derived from the OWR record, not an independently verified literature assessment. The supplied background expressly requests independent verification of time-sensitive status and cited claims.
\item Attribution in the supplied extract: question 9 was proposed by Trevor Wooley and attributed as a question of J.~Bober, D.~Fretwell, G.~Kopp, L.~Du, and T.~Wooley; question 10 was proposed by Trevor Wooley; question 11 by Jori Merikowski; question 12 by Ben Green, with its first part attributed to Erd\H{o}s and its second described as having been considered by Erd\H{o}s--Ruzsa and by Hildebrand; question 13 by Sarah Peluse and Robert Lemke Oliver; question 14 by Simon Myerson.
\item Historical assertions are attributed in the extract to Schinzel (1967) and Daniel Shiu (2000). No separate bibliographic references for these assertions are supplied. The report links above are supplied authoritative provenance, but their contents have not been visited or verified in preparing this draft.
\end{itemize}
\end{document}
